\chapter{Derivative problems that collapse}

\begin{orient}
\textbf{What this chapter is for.} There is a family of problems, common in competitions and common in examinations built by people who write competition problems, whose correct solution is to do almost nothing. They are engineered so that the intimidating surface hides a structure that makes the calculation trivial: the point at which the derivative is requested kills most of the terms, or the expression is identically constant, or the requested order of differentiation exceeds the degree available, or the outer layer of a composite cannot affect the answer. A solver who differentiates on reflex will get the right answer after two pages, or, more often, get a wrong answer after two pages. This chapter teaches the diagnostic that comes before differentiating, which is to evaluate the pieces at the requested point and ask which of them survive. It then teaches the opposite discipline, because the second-commonest failure here is inventing a collapse that is not there and reporting a clean wrong answer where the honest ugly one was correct.
\end{orient}

\section{The diagnostic}

Every technique in this chapter is a special case of one habit: read the whole question, including the evaluation point, before writing a single derivative. Differentiation is a local operation, and if the problem specifies where it wants the answer, that specification is information. Most of the time it is decisive information, because a problem that names a point has usually named it for a reason.

The diagnostic is three questions, asked in order, and it takes about twenty seconds.

\emph{First, is a derivative even being asked for?} A problem can be filed under differentiation and still be a geometry problem, a number theory problem, or an algebra problem with one trivial differentiation at the end. \emph{Second, does the expression simplify identically?} Baroque constructions built from trigonometric identities, conjugate logarithms, and geometric series are usually baroque on purpose, and the purpose is to hide a constant or a linear function. \emph{Third, if a specific point is named, what dies there?} Evaluate every factor and every inner function at that point before applying any rule, and mark the ones that vanish. In a product of $k$ factors the derivative has $k$ terms, and a single vanishing factor can eliminate $k-1$ of them.

\begin{center}
\begin{tikzpicture}[node distance=7mm and 12mm]
\node[dtest] (a) {Does the problem actually\\ ask for a derivative?};
\node[dleaf, right=12mm of a] (a1) {Answer the real question;\\ differentiate at the end, if at all};
\node[dtest, below=7mm of a] (b) {Is the expression identically\\ constant, or reducible?};
\node[dleaf, right=12mm of b] (b1) {Simplify to the end,\\ then differentiate what is left};
\node[dtest, below=7mm of b] (c) {Is one evaluation\\ point specified?};
\node[dleaf, right=12mm of c] (c1) {Evaluate every factor there;\\ delete the dead terms};
\node[dleaf, below=7mm of c] (c2) {Differentiate honestly;\\ strip inert wrappers first};
\draw[dedge] (a) -- node[dlab,above]{no} (a1);
\draw[dedge] (a) -- node[dlab,left]{yes} (b);
\draw[dedge] (b) -- node[dlab,above]{yes} (b1);
\draw[dedge] (b) -- node[dlab,left]{no} (c);
\draw[dedge] (c) -- node[dlab,above]{yes} (c1);
\draw[dedge] (c) -- node[dlab,left]{no} (c2);
\end{tikzpicture}
\end{center}

The tree is not a decoration. Three of its four leaves are cheap, and the expensive leaf at the bottom is reached only after the cheap ones have been ruled out. That ordering is the entire content of the chapter.

\section{Points that kill terms}

Start with the simplest and most frequent case, where the collapse is caused by the evaluation point alone and requires no cleverness at all to detect, only the willingness to look.

\tech{Evaluation point where most terms vanish}{medium}
\cue A product or quotient of three or more factors, with the derivative requested at one suspiciously convenient point: $x=0$, $x=1$, $x=\pi/2$, $x=\ln 2$.
\method Write the product rule symbolically before computing anything, as $f'=u'vw+uv'w+uvw'$, and then evaluate each factor at the point and delete every term containing a factor that is zero there. Note carefully which factor is differentiated in each term: if $u(a)=0$, the surviving term is the one containing $u'$, because that is the term in which $u$ itself does not appear. The standard annihilators are worth memorising. At $x=0$: $\sin x$, $\tan x$, $\arcsin x$, $\arctan x$, $\ln(1+x)$, $\eu^{x}-1$, and every positive power $x^{k}$. At $x=\pi/2$: $\cos x$ and $\cot x$. At $x=1$: $\ln x$ and $x-1$.

\begin{worked}
The classical shape:
\[
f(x)=\frac{\eu^{x^{2}}\sin x}{1+x^{3}},\qquad f'(0)=?
\]
At $x=0$ the numerator is $0$ and the denominator is $1$, so in the quotient rule $\dfrac{N'D-ND'}{D^{2}}$ the second term dies outright. Also $D=1$ and $D'=3x^{2}=0$ there, so $f'(0)=N'(0)$. And $N'=\eu^{x^{2}}(2x\sin x+\cos x)$, whose value at $0$ is $1$. Hence $f'(0)=1$, confirmed by a central difference as $1.0000000$. No quotient rule was ever executed.

A polynomial version where the collapse is total:
\[
f(x)=\prod_{k=1}^{10}(x-k),\qquad f'(1)=?
\]
The product rule gives ten terms, each a product of nine factors. Nine of those ten terms still contain the factor $(x-1)$, which vanishes at $x=1$. Only the term in which $(x-1)$ was differentiated survives, and it equals $\prod_{k=2}^{10}(1-k)=(-1)(-2)\cdots(-9)=(-1)^{9}9!=-362880$, confirmed numerically as $-362880.00$. Expanding the degree-ten polynomial first would be a page of arithmetic for the same number.

When two factors vanish, the first derivative is zero and the question moves up an order. For
\[
f(x)=\sin x\cdot\ln(1+x)\cdot\bigl(\eu^{x}+2\bigr),
\]
two of the three factors vanish at $0$, so $f'(0)=0$ without computation, and the counting rule gives $f''(0)=2!\cdot\sin'(0)\cdot\bigl(\ln(1+x)\bigr)'(0)\cdot\bigl(\eu^{0}+2\bigr)=2\cdot1\cdot1\cdot3=6$. A second difference returns $6.0000000$. Appending a fourth vanishing factor $\tan x$ pushes it to third order, and $f'''(0)=3!\cdot1\cdot1\cdot1\cdot3=18$, confirmed numerically as $18.00005$. Nothing here required the product rule to be written out at all.

And one where the annihilator is $\ln x$ at $x=1$:
\[
f(x)=\ln x\cdot\bigl(x^{3}+\arctan x\bigr)\cdot\eu^{x^{2}-1},\qquad
f'(1)=\frac{1}{x}\bigl(x^{3}+\arctan x\bigr)\eu^{x^{2}-1}\Big|_{1}=1+\frac{\pi}{4},
\]
since the other two terms both retain the factor $\ln x$. Numerically $1.7853982$ against $1+\pi/4=1.7853982$. The two remaining factors never had to be differentiated at all.
\end{worked}

\begin{trap}
Substituting the point into $f$ and then differentiating. That gives $0$ every time, and it gives $0$ for the wrong reason. The point is substituted into $f'$, after the symbolic differentiation, never before. The subtler error is deleting the wrong term: with $u(a)=0$ it is the terms containing $u$ undifferentiated that die, so the surviving term is the one carrying $u'$. Students who delete ``the term with the zero in it'' half the time delete precisely the term that survives.
\end{trap}

\begin{seealsob}The expression is secretly constant; the decoy evaluation point; the annihilated mixed partial.\end{seealsob}

\begin{keynote}[Counting the survivors before you compute them]
For $f=u_{1}u_{2}\cdots u_{k}$ the product rule has $k$ terms, and the $j$th term omits $u_{j}$ and includes $u_{j}'$. So the number of surviving terms at a point $a$ is decided by how many of the $u_{i}$ vanish at $a$. If exactly one vanishes, say $u_{m}(a)=0$, then exactly one term survives, the $m$th, and
\[
f'(a)=u_{m}'(a)\prod_{i\neq m}u_{i}(a).
\]
If two or more vanish, every term still contains at least one of them and $f'(a)=0$ outright, with no computation at all. That dichotomy, one zero factor or two, decides the entire problem before a single derivative is taken, and it extends: if $r$ factors vanish at $a$ then the first $r-1$ derivatives of $f$ vanish there and $f^{(r)}(a)=r!\prod$ of the surviving pieces times the product of the $u_{i}'(a)$ for the vanishing ones.
\end{keynote}

\begin{keynote}[Two more collapses that live at the origin]
The point $x=0$ carries more structure than any other, and two further species live only there. The first is the Maclaurin reading: a derivative requested \emph{only} at the origin is a series coefficient, since $f^{(n)}(0)=n!\,[x^{n}]f(x)$. For $f(x)=\arctan\bigl(x^{3}\bigr)$ the expansion is $x^{3}-\tfrac13x^{9}+\tfrac15x^{15}-\cdots$, so $f^{(9)}(0)=9!\cdot\bigl(-\tfrac13\bigr)=-120960$ and $f^{(10)}(0)=0$, and nobody differentiated anything nine times. The second is parity, which is a symmetry killing a term rather than a point killing one. If $f$ is even, every odd derivative vanishes at $0$; if $f$ is odd, every even one does. So for the even function $f(x)=\eu^{-x^{2}}\cos\bigl(x^{2}\bigr)$, $f^{(101)}(0)=0$ at no cost, and a problem that asks for exactly that is asking whether you noticed. Both species share this chapter's diagnostic: the point was named because the point does the work.
\end{keynote}

The vanishing-point collapse is local: it happens at the named point and nowhere else. The next kind is global. The expression is not merely convenient at one point, it is the same at every point, and every derivative beyond the first is zero.

\section{Expressions that are constant}

\tech{The expression is secretly constant}{medium}
\cue A baroque construction with a derivative requested at an arbitrary-looking point such as $x=\pi/5$ or $x=67$: nested trigonometry, an infinite geometric series whose ratio is a trigonometric square, a product of conjugate logarithms, or a functional inequality asserted for all real arguments.
\method Test the hypothesis ``this is identically constant'' before doing anything else. Two cheap tests: substitute two convenient values and compare, or simplify one layer at a time from the inside out. The engines that produce these are a short list. $\sin^{2}+\cos^{2}=1$ and $\sec^{2}-\tan^{2}=1$; the reflection identity $\arctan x+\arctan(1/x)=\pi/2$ for $x>0$; a product $\ln(a+b)\ln(a-b)$ arranged to conjugate; a geometric series whose sum is exactly the reciprocal of the prefactor; and a functional inequality valid for all arguments, where free choice of the arguments forces rigidity.

\begin{worked}
\[
f(x)=x\cdot\frac{1+\cos 2x}{2}\cdot\sum_{n\geq0}\left(\frac{\tan^{2}x}{1+\tan^{2}x}\right)^{n},
\qquad f'(\pi/5)=?
\]
Peel from the inside. The ratio is $\dfrac{\tan^{2}x}{\sec^{2}x}=\sin^{2}x$, so the series is geometric with ratio $\sin^{2}x<1$ and sums to $\dfrac{1}{1-\sin^{2}x}=\sec^{2}x$. The middle factor is $\cos^{2}x$ by the double-angle identity. So $f(x)=x\cos^{2}x\sec^{2}x=x$ identically, and $f'\equiv1$. In particular $f'(\pi/5)=1$. Evaluating the original expression numerically at $x=\pi/5$ returns $0.628318530718=\pi/5$ exactly, and the difference quotient returns $1.0000000$.

The reflection identity deserves its own line, because it is the engine most often disguised. For $x>0$, $\arctan x+\arctan\tfrac1x=\tfrac\pi2$, and for $x<0$ the same sum is $-\tfrac\pi2$. Either way the function is locally constant, so its derivative is $0$ everywhere except at $x=0$, where it is undefined. Differentiating directly confirms it and shows why the cancellation is exact:
\[
\dvop{x}\Bigl[\arctan x+\arctan\tfrac1x\Bigr]=\frac{1}{1+x^{2}}+\frac{1}{1+\tfrac1{x^{2}}}\cdot\Bigl(-\frac{1}{x^{2}}\Bigr)=\frac{1}{1+x^{2}}-\frac{1}{x^{2}+1}=0 .
\]
The useful lesson is in the constant: a problem that asks for the \emph{value} rather than the derivative needs the sign of $x$, whereas a problem that asks for the derivative does not. Reading which is wanted is the whole difficulty.

A second engine, shorter and even more common, is the conjugate logarithm. The pair
\[
f(x)=\ln\bigl(x+\sqrt{x^{2}+1}\bigr)+\ln\bigl(\sqrt{x^{2}+1}-x\bigr)
\]
looks like two unrelated inverse-hyperbolic pieces, but the two arguments are conjugates, and their product is $\bigl(\sqrt{x^{2}+1}\bigr)^{2}-x^{2}=1$. So $f(x)=\ln 1=0$ identically, and $f'\equiv0$ everywhere. Evaluated numerically at $x=0.3$, $2$ and $7$, the expression returns $5.6\times10^{-17}$, $4.4\times10^{-16}$ and $3.6\times10^{-15}$: zero, to floating-point precision. Differentiating the two logarithms separately is entirely feasible, produces $\dfrac{1}{\sqrt{x^{2}+1}}$ and $-\dfrac{1}{\sqrt{x^{2}+1}}$, and takes ten times as long to reach the same conclusion. The tell is always the same: two factors or two summands whose arguments differ by a sign in one place.

Now the functional-inequality version, which looks like a different subject. Suppose $f:\R\to\R$ satisfies
\[
f(x+y)+f(y+z)+f(z+x)\ \geq\ 3f(x+2y+3z)\quad\text{for all real }x,y,z .
\]
Find $f'(67)$. The substitution $u=x+y$, $v=y+z$, $w=z+x$ is invertible, since the coefficient matrix has determinant $2$, and $x+2y+3z=2v+w$. So the hypothesis says $f(u)+f(v)+f(w)\geq3f(2v+w)$ for all real $u,v,w$. Choose $v=t$ and $w=-t$, so that $2v+w=t$:
\[
f(u)+f(t)+f(-t)\geq3f(t)\ \Longrightarrow\ f(u)\geq 2f(t)-f(-t)\quad\text{for every }u .
\]
The right side does not involve $u$, so $f$ is bounded below; write $m=\inf f$, which is finite. Taking the infimum over $u$ gives $m\geq 2f(t)-f(-t)$, and replacing $t$ by $-t$ gives $m\geq 2f(-t)-f(t)$. Adding the two, $2m\geq f(t)+f(-t)$. But $f(t)\geq m$ and $f(-t)\geq m$, so $f(t)+f(-t)\geq 2m$. Both inequalities force equality, and since each summand is at least $m$, each equals $m$. Hence $f\equiv m$ and $f'(67)=0$.
\end{worked}

\begin{trap}
Differentiating the intimidating form correctly, obtaining a mess that ``should'' simplify, and then failing to make it simplify. The work is not wrong; it is wasted, and it is wasted in a way that is hard to abandon because each individual step was correct. The rule is to spend twenty seconds testing for constancy before spending twenty minutes not testing for it. The reverse trap is declaring constancy from two agreeing sample points: two points prove nothing, and the identity must actually be exhibited.
\end{trap}

\begin{seealsob}Evaluation point where most terms vanish; the decoy evaluation point; the outer wrapper is inert.\end{seealsob}

A constant function is the extreme case of a function with low degree. The next technique is the same collapse counted more finely: the expression is not constant, but it is of low degree in each variable separately, and the derivative requested is of high order in each.

\section{Orders that exceed the degree}

\tech{The annihilated mixed partial}{hard}
\cue A very high mixed partial of an expression assembled from $\ln x$, $\ln y$, and terms linear in one of the variables, often presented as a determinant or a sum of a dozen pieces, at an elaborate evaluation point.
\method Expand the expression into a sum of terms first, then test each term against the two annihilation rules. A term depending on only one variable is killed by any mixed partial. A term of degree at most $k-1$ in $x$ is killed by $\partial_{x}^{k}$, whatever else it contains, and a factor like $\ln y$ carries degree zero in $x$. If every term dies, the answer is $0$ and the ornate evaluation point was decoration.

\begin{worked}
\[
f(x,y)=\det\begin{pmatrix}\ln x&\ln y&1\\[1pt] \ln y&\ln x&1\\[1pt] x&y&1\end{pmatrix},
\qquad \frac{\partial^{110}f}{\partial x^{90}\,\partial y^{20}}\bigg|_{(67,69)}=?
\]
Expand the determinant once, along the bottom row or by the rule of Sarrus:
\[
f=(\ln x)^{2}-(\ln y)^{2}-x\ln x+x\ln y-y\ln x+y\ln y .
\]
Now classify the six terms. Four of them, namely $(\ln x)^{2}$, $-x\ln x$, $-(\ln y)^{2}$ and $y\ln y$, depend on one variable only, so any mixed partial annihilates them. The two cross terms survive that test but fail the degree test: $x\ln y$ is degree $1$ in $x$, and $\partial_{x}^{90}$ needs degree at least $90$; $-y\ln x$ is degree $1$ in $y$, and $\partial_{y}^{20}$ needs degree at least $20$. So
\[
\frac{\partial^{110}f}{\partial x^{90}\partial y^{20}}\equiv 0 ,
\]
identically, and the point $(67,69)$ was never relevant. Numerically, the fourth mixed difference $\partial_{x}^{2}\partial_{y}^{2}f$ at $(67,69)$, taken at step $0.5$, returns $-2\times10^{-13}$, which is round-off against zero.

The threshold matters and is worth checking. The same function has $\partial_{x}^{2}\partial_{y}f=1/x^{2}$, which at $x=67$ is $0.00022277$, against a numerical value of $0.00022279$. That derivative is not zero, because $\partial_{y}^{1}$ does not reach the degree threshold on the term $-y\ln x$. The collapse needs both orders to be at least $2$, and asserting it without checking both is guessing.
\end{worked}

\begin{trap}
Reaching for Jacobi's formula, $\dvop{t}\det A=\operatorname{tr}\bigl(\operatorname{adj}(A)A'\bigr)$, and applying it ninety times, instead of expanding one $3\times3$ determinant once. Any technique whose cost grows with the order of the derivative is the wrong technique when the order is in the dozens. The other error is asserting annihilation from the ``looks linear'' impression without expanding: an unexpanded determinant hides which terms are products of the two variables.
\end{trap}

\begin{seealsob}Mixed-partial coefficient extraction; the expression is secretly constant; parity annihilation.\end{seealsob}

The three collapses so far are triggered by different features of the problem, but they share a shape: some quantity in the expression is smaller than the operation applied to it. In the first, the number of nonvanishing factors at the point is smaller than the number of terms in the product rule. In the second, the number of essentially distinct values of the function is one. In the third, the degree in each variable is smaller than the order of differentiation in that variable. In each case the correct response is to establish the inequality and stop, and in each case the establishing takes a few lines while ignoring it takes a few pages. What follows is the discipline that keeps this reflex from becoming a superstition.

\section{The opposite error}

Everything so far trains one reflex: look for the collapse. That reflex, left unbalanced, produces its own characteristic failure. A solver who has learned that named points are significant will start manufacturing significance, will decide that $x=\pi-1$ must simplify because $\pi$ appeared, and will report a clean wrong answer with confidence. The correct posture is two-sided, and the discipline is symmetric: a collapse must be exhibited, not assumed, and an honest transcendental answer is a perfectly good answer.

\tech{The decoy point, and the large constant}{medium}
\cue An evaluation point that looks as though it should trigger a named-angle simplification but does not, such as $x=\pi-1$, $x=2$, $x=4$; or a large specific constant such as $2026$ that might be structural or might be noise.
\method Two-sided discipline. Going one way: if the point does not simplify anything, accept that, and execute the chain rule cleanly, reporting the transcendental value. A wrong clean answer is worse than a right ugly one. Going the other way: before dismissing a large constant as arbitrary, check whether it makes an identity close. Constants adjacent to a prime, or one less than the index of a root of unity, or one more than a factorial, are the ones to suspect. The test is cheap: try the same construction with the constant replaced by $3$ or $5$ and see whether a pattern appears.

\begin{worked}
\emph{The decoy.} Let $f(x)=\sin\bigl(\cos(1+x^{2})\bigr)$ and find $f'(\pi-1)$. The appearance of $\pi$ invites the hope that $1+x^{2}$ lands on a named angle. It does not: $1+(\pi-1)^{2}=5.5864191$, which is not a rational multiple of $\pi$. So differentiate honestly, outside in:
\[
f'(x)=-2x\cos\bigl(\cos(1+x^{2})\bigr)\sin\bigl(1+x^{2}\bigr),
\]
and $f'(\pi-1)=1.9791992$, matching the central difference $1.9791992$. That is the answer. There is nothing further to extract.

\emph{The structural constant.} Let
\[
f(x)=\prod_{k=1}^{2026}\Bigl(x^{2}-2x\cos\tfrac{k\pi}{2027}+1\Bigr),\qquad f'(1)=?
\]
Here $2026$ is not noise: it is one less than $2027$, and each factor is $\abs{x-\eu^{\iu k\pi/2027}}^{2}$ for a $4054$th root of unity. The identity is
\[
\prod_{k=1}^{n-1}\Bigl(x^{2}-2x\cos\tfrac{k\pi}{n}+1\Bigr)=\frac{x^{2n}-1}{x^{2}-1}=\sum_{j=0}^{n-1}x^{2j},
\]
checked numerically at $n=3,5,8,13$ and several values of $x$, all agreeing to ten digits. With $n=2027$ the product is $\sum_{j=0}^{2026}x^{2j}$, a polynomial, so
\[
f'(1)=\sum_{j=0}^{2026}2j=2\cdot\frac{2026\cdot2027}{2}=2026\cdot2027=4106702 .
\]
The general formula $f'(1)=n(n-1)$ was confirmed numerically at $n=3,5,13$, returning $6.000$, $20.000$ and $156.000$. Two thousand factors, and the derivative is a single multiplication.

\emph{The same constant, twice, once structural and once not.} Consider $f(x)=x^{2026}\eu^{-x}$ and the request for $f'(2026)$. Logarithmic differentiation gives $\dfrac{f'}{f}=\dfrac{2026}{x}-1$, which vanishes exactly at $x=2026$. So $f'(2026)=0$, and the constant in the exponent and the constant in the evaluation point were the same number for a reason: the point is the maximum of the integrand of $\Gamma(2027)$. Now consider $g(x)=\sin(2026x)$ and the request for $g'(1)$. Here $2026\cos(2026)=-1918.4724$, and $\cos(2026)=-0.9469262$ is not any recognisable number. The constant is decoration. The two problems look alike on the page and only the cheap test tells them apart.
\end{worked}

\begin{trap}
Manufacturing a simplification that is not there. If you cannot write down the identity you are appealing to, you do not have one. The mirror-image trap is dismissing a large constant without the cheap test: substituting $n=3$ or $n=5$ into the same construction takes thirty seconds and either exposes the pattern or rules it out.
\end{trap}

\begin{seealsob}Evaluation point where most terms vanish; the expression is secretly constant; the label is a decoy.\end{seealsob}

\begin{keynote}[Both errors have the same cure]
Assuming a collapse and missing a collapse are the same mistake made in opposite directions, and both are cured by the same twenty seconds of inspection. Evaluate the pieces. Write down the identity you think applies, in full. If it is there, the problem is finished; if it is not, you have lost twenty seconds and gained certainty that the ugly answer is the right one.
\end{keynote}

\subsection{What counts as evidence}

The two-sided discipline is only usable if you can say what evidence licenses a collapse, so here is the standard. An algebraic identity counts if you can write it down in full and check it at one numerical value. A vanishing factor counts if you have evaluated that factor at the point and got zero. A degree argument counts if you have expanded the expression into a sum of terms and inspected each one. A symmetry counts if you have named the map under which the expression is invariant.

What does not count: that the point contains $\pi$; that the constant is large and specific; that the expression is too ugly to be the intended answer; that two sample values happened to agree. The first three are the reasons students most often give themselves, and all three are compatible with there being no collapse whatever. The fourth is worse than useless, because a construction can be constant on a set of measure zero, and because two points also lie on a line.

It is worth keeping the five signatures in one list, because they are what you are scanning for and a list of five is short enough to run through in a few seconds. One: a named evaluation point at which a factor, an inner function, or a denominator's derivative vanishes. Two: an expression assembled from identity pairs, conjugates, or a summable series, so that the whole thing reduces to a constant or to $x$. Three: a differentiation order that exceeds the degree available in the relevant variable. Four: a monotone outer wrapper on a question about location or sign. Five: a statement whose real content is not calculus at all. Every problem in this chapter is one of those five, and most engineered problems are one of the first two.

There is one asymmetry worth exploiting. Checking for a collapse is cheap and bounded: evaluate the factors, sum the series, try $n=3$. Failing to check is unbounded, because the honest calculation of an engineered expression can run for pages. So the check is always worth making, even when you are fairly sure it will fail, and the only discipline required is to stop checking once it does.

\section{Layers that cannot matter}

The next collapse is not about the point or the constant but about the shape of a composite. Questions about \emph{where} a derivative vanishes, or about its \emph{sign}, are often insensitive to the outermost layer of the function, and recognising that saves the entire calculation.

\tech{The outer wrapper is inert}{hard}
\cue A question about stationary points, monotonicity, or the sign of a derivative for $g\bigl(h(x)\bigr)$, where $g$ is strictly monotone: $\exp$, $\ln$, $\arctan$, an odd power, or the Lambert $W$.
\method If $g'$ is nonzero throughout the relevant range, then $\bigl(g\circ h\bigr)'=g'(h)\,h'$ vanishes exactly where $h'$ vanishes and carries the sign of $h'$ wherever $g'>0$. Differentiate only the inside. The wrapper is inert for questions about location and sign, and it is emphatically \emph{not} inert for the value of the derivative or for the domain of the composite, so check what the question is asking before discarding it.

\begin{worked}
Let $f(x)=W\bigl(\sin\Gamma(x)\bigr)$ on $[2,\tfrac72]$, where $W$ is the principal Lambert function. Locate the stationary points.

The Lambert function satisfies $W'(u)=\dfrac{W(u)}{u\bigl(1+W(u)\bigr)}$ on the principal branch, an expression with its own branch analysis, and it is never zero. So the stationary points of $f$ are exactly those of $\sin\Gamma(x)$, whose derivative is $\cos\bigl(\Gamma(x)\bigr)\Gamma'(x)$. Now $\Gamma'=\Gamma\dig$, and the only positive zero of the digamma function is at $x=1.4616321$, which is outside $[2,\tfrac72]$; on that interval $\dig>0$, since $\dig(2)=0.4227843$, so $\Gamma'>0$ throughout and contributes no stationary point. What remains is $\cos\Gamma(x)=0$, that is $\Gamma(x)=\pi/2$, whose solution on the interval is
\[
x^{*}=2.7209149,\qquad \Gamma(x^{*})=1.5707963=\tfrac{\pi}{2}.
\]
That is the unique stationary point, and it is a maximum, with $f(x^{*})=W(1)=0.5671433$, the omega constant. Numerically $f'(2.6)=+0.0549981$ and $f'(2.85)=-0.0978350$, straddling it, and $f'(x^{*})=0$ to machine precision. At every sample point the sign of $f'$ agrees with the sign of $\bigl(\sin\Gamma\bigr)'$, as the argument predicts.

The same argument in a much more ordinary setting: to maximise $x^{1/x}$ for $x>0$, do not differentiate $x^{1/x}$. Since $\ln$ is strictly increasing, the maximiser of $x^{1/x}$ is the maximiser of $\dfrac{\ln x}{x}$, whose derivative is $\dfrac{1-\ln x}{x^{2}}$, zero at $x=\eu$. The maximum value is $\eu^{1/\eu}=1.4446679$. Differentiating $x^{1/x}$ directly is possible, needs logarithmic differentiation anyway, and produces the factor $x^{1/x}$ that was never going to vanish. Every optimisation of a product, a quotient, or a power is an instance of this: take logarithms, because the logarithm is the archetypal inert wrapper.

Now the part the inertness argument does \emph{not} cover. Extend the question to $[2,4]$. The principal $W$ requires its argument to be at least $-1/\eu=-0.3678794$, and $\sin\Gamma(x)$ drops below that: it crosses $-1/\eu$ downward at $x=3.5512873$ and returns at $x=3.9874729$, reaching $-1$ at $x=3.8032322$, where $\Gamma(x)=3\pi/2$. So $f$ is undefined on an interval of length $0.436$, and the stationary point of $\sin\Gamma$ at $3.8032322$ is not a stationary point of $f$ because $f$ does not exist there. The wrapper cannot move a stationary point, but it can delete one by removing it from the domain.
\end{worked}

\begin{trap}
Computing $W'$, or $\Gamma'$ in closed form, only to find that neither vanishes. The cost is high and the information gained is zero. The opposite trap is the one exposed above: treating the wrapper as inert for a question it does not govern. ``Where is $f'=0$'' and ``what is the sign of $f'$'' are insensitive to a monotone outer layer; ``what is $f'(a)$'', ``where is $f$ defined'', and ``how many stationary points does $f$ have'' are not.
\end{trap}

\begin{seealsob}Sign charts and extremum counting; chain-rule discipline; the expression is secretly constant.\end{seealsob}

\section{When the label lies}

The last collapse is the most complete. Some problems filed under differentiation contain no function to differentiate, or contain one whose differentiation is a single line arriving after all the real work is done. The signature is that the statement is long and the word ``derivative'' appears once, near the end.

\tech{The label is a decoy: no derivative is required}{hard}
\cue A problem headed ``derivative'' whose statement is mostly geometry, number theory, or algebra, and in which the object to be differentiated is not even named until the final clause; or a ratio of two evaluations of the same unknown polynomial.
\method Read the question, not the heading. Identify what is actually unknown and what the constraints actually determine. The recurring cores hiding under a derivative label are worth listing: Vieta relations on a quadratic whose roots are trigonometric functions of the angles of a triangle; complex-region geometry reduced to a point-to-line distance; a divisor count $\tau(N)$ feeding into $\Gamma$ of a small integer; the reflection symmetry of $\ln x$ and $\eu^{x}$ about the line $y=x$; and a fixed point of an iterated map that turns a tower of exponents into a constant. In every case the differentiation, when it comes, is one line.

\begin{worked}
\emph{Vieta under a derivative heading.} In right triangle $PQR$ with $\angle R=\pi/2$, the numbers $\tan\tfrac P2$ and $\tan\tfrac Q2$ are the roots of $Q(x)=a(x^{2}-sx+p)$. Compute $\dfrac{Q'(1/2)}{Q(0)}$.

Nothing about this requires differentiating anything interesting. Since $P+Q=\pi/2$, we have $\tfrac{P+Q}{2}=\tfrac\pi4$, and the tangent addition formula gives
\[
\tan\frac{P+Q}{2}=\frac{\tan\tfrac P2+\tan\tfrac Q2}{1-\tan\tfrac P2\tan\tfrac Q2}=\frac{s}{1-p}=1,
\]
so $s=1-p$, that is $s+p=1$. Now the one line of calculus: $Q'(x)=a(2x-s)$, so $Q'(1/2)=a(1-s)=ap$. And $Q(0)=ap$. Hence the ratio is $1$, independently of $a$ and of the triangle. Random right triangles and random $a$ reproduce $s+p=1.0000000$ and the ratio $1.0000000$.

\emph{Geometry under a derivative heading.} Find the minimum distance between the curves $y=\ln x$ and $y=\eu^{x}$. The tempting move is to set up the squared distance between $(x,\eu^{x})$ and $(t,\ln t)$ and solve a two-variable stationarity system. Do not. The curves are reflections of one another in the line $y=x$, so the segment realising the minimum is perpendicular to that line, and the problem reduces to minimising the distance from $y=\eu^{x}$ to the line, namely $\dfrac{\eu^{x}-x}{\sqrt2}$. One derivative: $\eu^{x}-1=0$ at $x=0$, giving distance $\dfrac1{\sqrt2}$ from the curve to the line. Doubling, the distance between the curves is $\sqrt2$, attained between $(0,1)$ and $(1,0)$. A brute-force descent over both curves converges to $1.4142136$ with closest points $(0,1)$ and $(1,0)$, against $\sqrt2=1.4142136$.

\emph{A fixed point under a derivative heading.} Let $y=x^{x^{x^{\cdots}}}$ and suppose $y=2$. Find $\dv{y}{x}$ there. The infinite tower is the fixed point of $u\mapsto x^{u}$, so it satisfies $y=x^{y}$, and this is the whole content of the tower notation: the exponent becomes the value. With $y=2$ the equation reads $2=x^{2}$, so $x=\sqrt2$. Only now does any calculus happen. Taking logarithms of $y=x^{y}$ gives $\ln y=y\ln x$, and differentiating,
\[
\frac{y'}{y}=y'\ln x+\frac{y}{x}\ \Longrightarrow\ y'=\frac{y^{2}}{x\bigl(1-y\ln x\bigr)} .
\]
At $x=\sqrt2$, $y=2$: $y\ln x=2\cdot\tfrac12\ln2=\ln2$, so $y'=\dfrac{4}{\sqrt2\,(1-\ln2)}=9.2175367$. Solving $y=x^{y}$ numerically on either side of $\sqrt2$ and differencing returns $9.2175367$. Iterating the tower at $x=\sqrt2$ converges to $2.0000000$, confirming that the branch chosen is the right one.
\end{worked}

\begin{trap}
Spending the whole time budget hunting for a derivative technique when the intended skill is elsewhere. The tell is that the statement contains a great deal of structure that no differentiation would use. A second and more specific error: when the extremum sits on an open boundary, an infimum exists but a minimum does not, and reporting a minimum is then wrong even though the number is right. Check that the optimising point is in the domain before naming it.
\end{trap}

\begin{seealsob}Constrained optimisation; the expression is secretly constant; the self-similar fixed point.\end{seealsob}

\section{Why these problems are worth studying}

It is tempting to file this chapter under recreation: engineered problems, competition tricks, curiosities that will never appear in an engineering calculation. That would be a mistake, and the reason is methodological. An ordinary exercise requires recognition and computation together, and shares its marks between them, so a solver who finishes one never learns which of the two was carrying the answer. The problems here are built so that the computation is unaffordable and only recognition can pay. Each of them is therefore a controlled experiment with a single factor varied, and the thing it measures is whether you looked before you started.

The structures being detected are not artificial, which is why the training transfers. A factor that vanishes at the point of interest is what a boundary condition looks like. An expression that is secretly constant is what a conservation law looks like. A differentiation order that exceeds the available degree is what a truncated model looks like. A monotone outer wrapper is what a change of units, a link function, or a log-likelihood looks like. No setter invented any of these; a setter merely arranges for one of them to be decisive and for every other route to be a wall. Meeting them here, where a shortcut is guaranteed to exist, is cheap practice for meeting them where nothing guarantees it and where nobody tells you afterwards that one was available.

The thesis of this chapter is narrow and worth stating plainly. Differentiation is cheap in principle and expensive in practice, and the expense is almost entirely in the terms you did not need. Every technique here removes terms before they are written: the vanishing point removes all but one of them, constancy removes all of them, insufficient degree removes all of them again, the monotone wrapper removes an entire layer, and the mislabelled problem removes the differentiation altogether. The habit that produces all five is the same, and it is the habit this whole book is about. Look at the object before you apply a rule to it, and let the structure decide the method rather than the other way round.

\section{Recognition matrix}
\begin{recog}{@{}p{0.31\textwidth}p{0.35\textwidth}p{0.28\textwidth}@{}}
\rowcolor{mist}\mh{If you see} & \mh{Reach for} & \mh{If that fails} \\
\midrule
\endhead
Three or more factors, derivative at one point & Write the product rule, then delete terms with a vanishing factor & Take logarithms and differentiate \\
$x=0$ in the statement & $\sin$, $\tan$, $\arcsin$, $\arctan$, $\ln(1+x)$, $x^{k}$ all vanish & Expand to first order in $x$ \\
$x=1$ with a $\ln x$ factor present & Only the term carrying $1/x$ survives & Substitute $x=1+h$ and expand \\
$\prod_{k}(x-a_{k})$, derivative at some $a_{j}$ & $\prod_{k\neq j}(a_{j}-a_{k})$; one term survives & Logarithmic differentiation \\
$f^{(n)}(0)$ for large $n$, only at the origin & Maclaurin coefficient: $n!\,[x^{n}]f$ & Expand the outer function in its own variable \\
$f$ even or odd, a derivative at $0$ & Parity kills the opposite half outright & Split $f$ into even and odd parts \\
Nested trig over a geometric series & Sum the series; suspect an identity & Evaluate at two points and compare \\
Conjugate logs, or $\arctan x+\arctan\tfrac1x$ & Constant; the derivative is $0$ & Simplify one layer at a time \\
A functional inequality for all arguments & Free substitution; expect rigidity, then $f'=0$ & Set the arguments equal, then opposite \\
$\dfrac{\partial^{n+m}f}{\partial x^{n}\partial y^{m}}$ with $n,m$ huge & Expand first; kill one-variable terms and low degrees & Coefficient extraction, $n!\,m!\,[x^{n}y^{m}]f$ \\
A determinant with a high mixed partial & Expand the determinant once, then classify terms & Check both orders exceed both degrees \\
A point like $\pi-1$ or $2$ that almost simplifies & Nothing; differentiate honestly & Report the transcendental value \\
A large constant such as $2026$ & Test whether $n-1$ completes a root-of-unity identity & Rerun the construction with $n=3$ \\
$W(\cdot)$, $\ln(\cdot)$, $\arctan(\cdot)$ wrapped round the real content & Stationary points and signs come from the inside only & Check the wrapper's domain separately \\
``Derivative'' but no function in the statement & Solve the real problem; differentiate at the end & Look for Vieta, symmetry, or a fixed point \\
Two curves that are mirror images in $y=x$ & Distance to the line, then double & Check the contact point lies in the domain \\
\end{recog}
